\documentclass[11pt]{article}
\usepackage[margin=.81in]{geometry}
\usepackage{amsmath,amssymb,booktabs,hyperref}
\hypersetup{colorlinks=true,urlcolor=blue}
\newcommand{\F}{\mathbb F}
\renewcommand{\labelitemi}{\(\bullet\)}
\title{\vspace{-1.6em}Exact \(p\)-adic root counts for six terms}
\author{A supplement to the ProveIt GRH article}
\date{27 November 2026}
\begin{document}
\maketitle
Let \(p\ge7\) be prime and \(v_p(p)=1\). For any nonzero center \(A\)
and depth \(r>1/(p-1)\), the exact maximum number of roots in the
closed disk \(v_p(X/A-1)\ge r\) of a nonzero polynomial with at most
six terms is
\begin{center}
\begin{tabular}{cc}
\toprule
Depth \(r\) & Maximum\\
\midrule
\(1/(p-1)<r\le1/(p-2)\) & \(2p\)\\
\(1/(p-2)<r\le1/(p-4)\) & \(p+1\)\\
\(1/(p-4)<r\le1/(p-5)\) & \(p\)\\
\(1/(p-5)<r\) & \(5\)\\
\bottomrule
\end{tabular}
\end{center}
The upper bound counts multiplicity. All extrema are attained by
distinct roots in finite extensions, with rational witnesses
about \(1\). They can also be attained by irreducible polynomials
over \(\mathbb Q\) about \(1\). At \(p=7\) this gives \(14,8,7,5\).
Neither RH nor GRH is used.

\section{What a minimal coefficient can see}
Use the \href{https://github.com/VladimirReshetnikov/ProveIt/blob/dd990b27a22290c756c1bd4adce1d971bf7e7821/Analysis/GeneralizedRiemannHypothesis/Article/sections/smd.tex#L57}{article's local transformation}
\(X=A\exp(bY)\) with \(v_p(b)=r\) (first take \(r\) rational).
The \(k\)-th coefficient is \(M_k b^k/k!\), where
\(M_k=\sum_j c_j A^{n_j}n_j^k\). The first six moments are not
all zero by Vandermonde; pad with unused exponents if necessary.
Normalize by a moment of least valuation among those first six.
The integral monic recurrence makes all moments integral, defining a
functional \(L\) on the quotient by \(\prod_j(T-n_j)\).

Write \(w(k)=rk-v_p(k!)=(r-1/(p-1))k+S_p(k)/(p-1)\).
A minimizing index \(k\ge6\) must have \(S_p(k)\le4\):
otherwise \(w(k)>w(5)\). Its moment has valuation less than one
because \(w(i)-w(k)<1\) for \(i\le5<k\).
If \(T^k\) is spanned modulo \(p\) by powers of strictly smaller
weight, then some such earlier-in-weight moment has valuation no
larger than \(M_k\). Hence \(k\) cannot minimize. This comparison
works in ramified fields since lifting the residue identity gives
an error divisible by \(p\).

In the residue algebra put \(U=T^p,\ V=T-U\). On the factor for
residue \(a\), \(U=a,\ V=T-a\), since the multiplicity is at most
six. Thus \(T^{p^j}=U\) for \(j\ge1\). A base-\(p\) digit expansion
compresses to \(T^iU^j\), \(i+j=S_p(k)\), where \(i\) is the units
digit and \(j\) the sum of higher digits; strict compression lowers
the weight. It remains to consider \(i+j\le4\).
\newpage
\section{The residue span calculation}
Apart from pure powers of \(T\) up to degree five, the only candidates
needed are \(U,TU,T^2U,U^2\). The last two are needed only for the
profile \((3,2,1)\) of residue multiplicities. Here are the reductions:
\begin{itemize}
\item With at most two residues \(U^2\) is affine in \(U\).
For \((U-a)T^2\), use \((U-a)(T-b)^2=0\) if \(b\)
has multiplicity at most two. In profile \((3,3)\) the expression
\(T^2U-T^3/3-(a+b)TU+abT\) has first and second derivatives zero
at both residues, so is affine in \(U\). For \(T^3U\), use
\((U-a)(T-b)^3=0\), with \(b\) of multiplicity at most three.
The one-residue case has \(U\) constant.
\item If \(a\) is the only repeated residue, then
\(TU=aT+U^2-aU\). For \(d=3,4\), reduce \(U^d\) by the polynomial
of the distinct residues if there are at most \(d\) of them.
Otherwise the multiplicity at \(a\) is at most \(d\), so
\((T-a)^d=(U-a)^d\). This again reduces \(U^d\).
For mixed terms use \(T^iU^j=U^{i+j}+a^j(T^i-U^i)\).
If there are no repeated residues then \(U=T\).
\item In the other profiles with \(V^2=0\), use \(U^2=2TU-T^2\).
For \(d=3,4\), first reduce \(U^d\) if there are at most \(d\)
residues; then use \(T^d=U^d+dU^{d-1}V\) to reduce the mixed powers.
The remaining case is \(d=3\) in \((2,2,1,1)\), where \(U^2V\)
is spanned by \(V,UV\) because \(V\) is supported on two distinct
residues. Reduce \(U^3\) using \(T^3\).
\item In \((3,2,1)\), the six classes \(1,T,T^2,U,TU,U^2\)
are a basis. Indeed, their span is the span of \(1,U,U^2,\allowbreak V,UV,V^2\);
values at three residues, first jets at two, and a second jet at
one identify these six classes independently. The basis handles
\(TU^2,U^3\) and all mixed candidates of total degree four.
\end{itemize}
Each reduction uses strictly earlier indices of no greater digit
sum, so these indices have strictly smaller weight.
The largest remaining index is \(2p\).

In \((3,2,1)\) let \(J\) vanish on the first five basis entries
and take value \(1\) on \(U^2\). For triple, double, and singleton
residues \(a,b,c\), respectively, \(J(T^3)=a+c-2b\).
One way to compute this is to use
\(J(V)=0,\ J(UV)=-1,\ J(V^2)=1\). Then
\(J(U^3)=a+b+c,\ J(U^2V)=-(a+b),\ J(UV^2)=a\).
If \(a+c-2b=0\), an affine change of variable gives \((0,1,2)\),
where \(J(T^4)=7-4=3\) using the same identities.
Thus \(U^2\) may be replaced by \(T^3\) or \(T^4\).
Beyond \(r=1/(p-2)\) every weight of that modified basis is
less than both \(w(p+2)\) and \(w(2p)\): in particular
\(w(p+2)-w(4)=(p-2)r-1\) and
\(w(2p)-w(4)=2((p-2)r-1)\).
Beyond \(1/(p-4)\), \(w(p+1)>w(5)\), and beyond \(1/(p-5)\),
\(w(p)>w(5)\). Apply Strassmann to obtain the upper bounds.
For a nonrational \(r\), choose a slightly smaller rational depth
in the same region for the upper bound.
\newpage
\section{Matching witnesses}
For the first endpoint \(r=1/(p-2)\), use the distinct exponents
\(0,p,2p,1,p+1,2\). In the integral exponent quotient the six classes
\(1,T,T^2,T^p,T^{p+1},T^{2p}\) are a basis over \(\mathbb Z_{(p)}\)
by their reduction just proved. Prescribe \(L=(0,0,0,0,0,1)\) on
this basis. Over \(\mathbb Q\) evaluations at the six exponents
span the dual, so \(L\) gives rational coefficients \(c_j\).
Its residue moments on \(T^3,T^4,T^{p+2}\) are \(0,3,1\),
as in the preceding calculation for \((0,1,2)\), including
\(J(T^2U)=J(U^3)+2J(U^2V)+J(UV^2)=1\).
At the endpoint the minimum analytic coefficient valuation is
\(4r\), attained exactly at \(4,p+2,2p\). To see this, use the
remaining-index list, the exact vanishing at \(0,1,2,p,p+1\),
and the divisibility by \(p\) at 3.
Weierstrass preparation in \(Y\) gives \(2p\) roots in the closed
disk counted with multiplicity.

For the endpoint \(r=1/(p-4)\), use \(0,p,2p,1,2,c\), where
\(c\) is the least positive residue of \(-2/3\bmod p\).
Take \(L\) to extract the \(T^5\)-coefficient in the exponent
quotient. In its residue algebra,
\[
 U = \beta T^4-\beta(3+c)T^3,\qquad \beta=(2c)^{-1},
\]
because the right side minus \(T\) is
\(\beta T(T-1)(T-2)(T-c)\). The right side agrees with \(T^p\)
on all four residue factors: it has zero coefficients in degrees \(0,1,2\),
and its value is \(T\) at \(1,2,c\).
It follows that \(M_p\) is divisible by \(p\), whereas
\(M_{p+1}\) and \(M_5\) are units. All moments at indices below
5 are exactly zero. The minimizing indices are \(5,p+1\):
the sole other candidate \(p\) satisfies \(1+w(p)>w(5)\).
Preparation gives \(p+1\) roots.

For \(r=1/(p-5)\), use
\[
 X^p-\sum_{k=0}^4\binom pk(X-1)^k+p^{11}.
\]
Its Newton polygon in \(X-1\) has vertices \((0,11),(5,1),(p,0)\).
All \(p\) roots are in the closed disk. For the deepest range take
five distinct rational roots sufficiently close to 1.

For the first three constructions, clear denominators, then keep
the leading term and perturb the other coefficients within their
support by a sufficiently high \(p\)-power. This preserves the
minimizing indices at the endpoint: the perturbation in the
analytic coefficients has valuation at least the chosen precision
plus \(w(k)\ge0\), or use the polynomial Newton polygon for the
third witness. By CRT choose the perturbation to be Eisenstein at
some other prime not dividing the leading coefficient.
This gives irreducibility and hence distinct roots. The same
perturbation works for the last construction in disjoint small
disks about the original rational roots. Endpoint witnesses also
work at all the shallower depths in their respective regions.

\paragraph{Prior art.}
The general order-\(n\) estimate \((n-2)p\) over \(\mathbb Q_p\)
for \(p>n\ge3\) is classical; see
\href{https://doi.org/10.1112/jlms/s2-36.1.1}{van der Poorten--Rumely}.
The result here specifies the exact depth-dependent profile in the
six-term, integer-exponent setting. No priority or Lean claim is made.
\newpage
\section{Partial certificate at discriminant \(-87\)}
This appendix concerns a different question; the general disk
profile does not settle it. In the ring
\(\mathbb Z_{(2)}[X]/(H_{-87})\), write \(\alpha\) for the class
of \(X\). The
\href{https://github.com/VladimirReshetnikov/ProveIt/blob/10802a6b36/Analysis/GeneralizedRiemannHypothesis/Research/grh_sparse_singular_moduli/checks/five_small_witnesses.json}{pinned class polynomial}
gives the exact identity modulo 32
\[
 \alpha^{63}\equiv1+16(\alpha+\alpha^2+\alpha^3).
\]
Reduction modulo 2 is the quotient by \(X^6+X^5+X^2+X+1\),
and the class of \(X\) has order 63. Write \(W=(\alpha^{63}-1)/16\)
and \(\ell=X+X^2+X^3\) for its reduction.

Given six distinct exponents, translate so they are nonnegative.
Group them as \(a+63q\) with \(0\le a<63\).
In the block with residue \(a\), take divided differences in \(q\)
of the columns \(\alpha^a(1+16W)^q\), of orders \(j=0,\dots,m-1\)
for multiplicity \(m\). This is invertible over \(\mathbb Q_2\).
Multiply the order-\(j\) column by \(j!/16^j\); the reduction
of the resulting integral column is \(X^a\ell^j\).
Indeed the \(j\)-th divided difference of \(\binom{q}{k}\) vanishes
for \(k<j\), equals \(1/j!\) when \(k=j\), and belongs to
\(k!^{-1}\mathbb Z_{(2)}\) when \(k>j\). The last statement follows
by taking the divided difference of the integer polynomial
\(k!\binom{q}{k}\) at integer nodes, for example by the monomial
formula in complete homogeneous symmetric polynomials.
The remaining powers of 16 make the higher terms even:
for \(k>j,\ 0\le j\le5\),
\(v_2(k!/j!)<4(k-j)\). This follows from the factorial digit
formula for \(k-j\ge2\) and \(v_2(j+1)\le2\) for \(k-j=1\).

Thus a full-rank residue matrix of columns \(X^a\ell^j\)
certifies \emph{all} lifts of that residue multiset. For a single
residue the six powers of \(\ell\) have bit masks
\(1,14,51,7,42,48\) and are independent, recovering the article's
single-class result. Across all profiles, exhaustively enumerating
weakly increasing multisets with at least one zero verifies
\(3{,}892{,}897\) out of \(\binom{67}{5}=9{,}657{,}648\).
Any pattern can be anchored at zero by a common exponent
translation. A failed rank test is inconclusive, and the full
case of \(-87\) remains open here.

\paragraph{Replay.}
The standard-library Python scripts check the six-slot jets,
the rational witness moments, the congruence modulo 32, and
the anchored multiset count. The written proof, rather than
finite testing, establishes the all-prime result.
\end{document}
